\renewcommand{\chaptername}{بخش}
\chapter{آزمایش‌ها و نتایج تجربی}

\section{داده‌ها و پیش‌پردازش}
در این بخش از مجموعهٔ دادهٔ dSprites استفاده شده است که شامل $64\times64$ تصاویر باینری و عوامل مولد حقیقی شناخته‌شده می‌باشد (شکل، مقیاس، زاویه، موقعیت‌های X و Y). پیش‌پردازش شامل نرمال‌سازی پیکسل‌ها به بازهٔ $[0,1]$، جداسازی تصادفی به نسبت 80/20 برای آموزش/اعتبارسنجی و عدم استفاده از افزایش داده به‌منظور حفظ عوامل مولد است.

\section{پیکربندی‌های آموزشی}
\subsection{سخت‌افزار و نرم‌افزار}
\begin{itemize}
  \item سخت‌افزار: GPU مبتنی بر CUDA یا Apple Silicon (MPS)
  \item چارچوب: PyTorch 2.0+
  \item نسخهٔ پایتون: 3.8+
  \item کتابخانه‌های کمکی: NumPy, Matplotlib, scikit-learn, NetworkX
\end{itemize}

\subsection{تنظیمات عمومی آموزش}
\begin{itemize}
  \item نرمال‌سازی تصاویر: مقادیر بین 0 و 1
  \item تقسیم داده‌ها: 80\% آموزش، 20\% اعتبارسنجی
  \item بهینه‌ساز: Adam (lr=1e-3)، batch size = 64
  \item تعداد epoch: 50 (یا تا همگرایی)
  \item جدول پیکربندی‌ها (نمونه):
\end{itemize}

\begin{table}[H]
\centering
\begin{tabular}{@{}lllll@{}}
\toprule
\textbf{مدل} & \textbf{$\beta$} & \textbf{ابعاد پنهان} & \textbf{ابعاد پنهان (شبکه)} & \textbf{اپوک} \\
\midrule
VAE-baseline & 1.0 & 10 & 256 & 50 \\
$\beta$-VAE-low & 2.0 & 10 & 256 & 50 \\
$\beta$-VAE-med & 5.0 & 10 & 256 & 50 \\
$\beta$-VAE-high & 10.0 & 10 & 256 & 50 \\
VAE-large & 1.0 & 20 & 512 & 50 \\
VAE-deep & 1.0 & 10 & 512 & 50 \\
\bottomrule
\end{tabular}
\caption{پیکربندی‌های آزمایشی}
\label{tab:exp_configs}
\end{table}

\section{شاخص‌های ارزیابی و روش‌شناسی}
ارزیابی شامل معیارهای کمی و کیفی است:
\begin{itemize}
  \item کیفیت بازسازی: MSE یا Binary Cross-Entropy
  \item تفکیک‌پذیری (Disentanglement): MIG، SAP، DCI
  \item تحلیل بصری: بازسازی، نمونه‌سازی و تراورس‌های فضای پنهان
\end{itemize}

برای محاسبهٔ MIG از روش مبتنی بر ماتریس اطلاعات متقابل بین ابعاد پنهان و عوامل استفاده شده است؛ سپس میانگین اختلاف بین بالاترین و دومین MI نرمال‌شده گرفته شده است.

\section{نتایج تجربی خلاصه}
جدول زیر نتایج جامع را نشان می‌دهد (مقادیر نمونه‌ای از گزارش اصلی):

\begin{table}[H]
\centering
\begin{tabular}{@{}lcccccc@{}}
\toprule
\textbf{مدل} & \textbf{$\beta$} & \textbf{MIG} & \textbf{Recon Loss} & \textbf{KL Loss} & \textbf{Total Loss} \\
\midrule
VAE & 1.0 & 0.15 & 0.082 & 4.52 & 4.602 \\
$\beta$-VAE (2) & 2.0 & 0.28 & 0.095 & 3.87 & 7.835 \\
$\beta$-VAE (5) & 5.0 & 0.42 & 0.134 & 2.91 & 14.684 \\
$\beta$-VAE (10) & 10.0 & 0.51 & 0.201 & 2.15 & 21.701 \\
VAE-large & 1.0 & 0.18 & 0.076 & 5.12 & 5.196 \\
VAE-deep & 1.0 & 0.16 & 0.079 & 4.88 & 4.959 \\
\bottomrule
\end{tabular}
\caption{خلاصهٔ نتایج آزمایشی (مقادیر نمونه‌ای)}
\label{tab:exp_summary}
\end{table}

\section{تحلیل و تفسیر}
خلاصهٔ نکات مهم:
\begin{itemize}
  \item امتیازات MIG با افزایش $\beta$ بهبود می‌یابد؛ اما خطای بازسازی افزایش می‌یابد (trade-off مشخص).
  \item ابعاد فعال فضای پنهان کاهش می‌یابد که نشان‌دهندهٔ فشرده‌سازی قوی‌تر است ولی احتمال فروپاشی پستریری (posterior collapse) را افزایش می‌دهد.
  \item برخی عوامل مانند شکل (discrete) راحت‌تر جدا می‌شوند نسبت به عوامل موقعیت (continuous).
\end{itemize}

\section{آزمایش‌های تکمیلی (Ablation)}
مطالعات ablation نشان می‌دهد که مواردی مانند KL annealing، scheduler بهبود همگرایی و استفاده از clipping برای پایداری اثرگذار هستند. نتیجه‌گیری‌ها مطابق با جدول نمونهٔ زیر است:

\begin{table}[H]
\centering
\begin{tabular}{@{}lccc@{}}
\toprule
\textbf{Ablation} & \textbf{MIG} & \textbf{Recon Loss} & \textbf{تغییر}\
\midrule
Full model ($\beta$=5) & 0.42 & 0.134 & Baseline \\
No KL annealing & 0.38 & 0.142 & -9.5\% MIG \\
No scheduler & 0.39 & 0.148 & -7.1\% MIG \\
Smaller encoder & 0.35 & 0.156 & -16.7\% MIG \\
No gradient clip & 0.40 & 0.151 & -4.8\% MIG \\
\bottomrule
\end{tabular}
\caption{نتایج مطالعات ablation (نمونه‌ای)}
\label{tab:ablation_summary}
\end{table}

\section{عمومیت‌پذیری بین-دیتاستی}
نتایج نشان می‌دهد که مدل‌های تمرین‌شده روی dSprites تا حدی به مجموعه‌های مشابه تعمیم می‌یابند اما برای داده‌های پیچیده‌تر (مانند SmallNORB) افت عملکرد قابل توجهی مشاهده می‌شود.

\section{خلاصهٔ بخش تجربی}
این بخش مجموعه‌ای از تنظیمات، نتایج کمی و تحلیل‌های کیفی را ارائه کرد. اعداد نمونه‌ای بر اساس نتایج گزارش اصلی هستند و در آماده‌سازی نهایی باید با خروجی‌های واقعی اجرا جایگزین شوند.

